\documentclass[12pt,a4paper,oneside]{extreport}
\usepackage{ctu-report}
\usepackage{amsmath}
\usepackage{placeins}
\usepackage{needspace}
\usetikzlibrary{positioning}
\setlength{\emergencystretch}{3em}
\renewcommand{\lstlistingname}{Đoạn mã}
\AtBeginDocument{\renewcommand{\thelstlisting}{\arabic{chapter}.\arabic{lstlisting}}}
\graphicspath{{assets/figures/}}
\Urlmuskip=0mu plus 1mu
\def\UrlBreaks{\do\/\do-\do.\do_}

% Figure helper: \ctufigure{file}{width as a fraction of \textwidth}{caption}{label}
% Height is capped so tall screenshots still fit on a page with their caption.
\newcommand{\ctufigure}[4]{%
  \begin{figure}[!htbp]
    \centering
    \includegraphics[width=#2\textwidth,height=0.62\textheight,keepaspectratio]{#1}
    \caption{#3}
    \label{#4}
  \end{figure}
  \FloatBarrier
}

% Caption for the unnumbered tables in the front matter and introduction.
\newcommand{\ctuplaintablecaption}[1]{%
  \begin{center}
    \itshape Bảng: #1
  \end{center}
  \addcontentsline{lot}{table}{Bảng: #1}%
}

% -----------------------------------------------------------------------------
% METADATA
% -----------------------------------------------------------------------------
\CTUReportType{NIÊN LUẬN}
\CTUTitle{Xây dựng quy trình CI/CD hoàn chỉnh cho ứng dụng Web}
\CTUStudent{Trần Thanh Duy}
\CTUStudentID{DC24V7X604}
\CTUCohort{2024--2027}
\CTUSupervisor{ThS. Nguyễn Minh Trung}
\CTUMonthYear{09/2026}

\begin{document}

\ctumakecover
\ctumakeinnertitle

% -----------------------------------------------------------------------------
% FRONT MATTER
% -----------------------------------------------------------------------------
\ctufrontmatter

\ctuheading{LỜI CẢM ƠN}
Em xin gửi lời cảm ơn chân thành đến các thầy cô tại Đại học Cần Thơ, đặc biệt
là các thầy cô trong trường Công Nghệ Thông Tin và Truyền Thông. Trong hơn một
năm qua các thầy cô đã tận tình giảng dạy, hướng dẫn nhiệt tình để em có được
những kiến thức mới đầy hữu ích cho tương lai sau này.

Sau gần 3 tháng tìm hiểu và thực hiện, đề tài ``Xây dựng quy trình CI/CD hoàn
chỉnh cho ứng dụng Web'' đã hoàn thành, ngoài sự nỗ lực của bản thân, em còn
nhận được nhiều sự động viên, khích lệ từ thầy cô và bạn bè.

Em xin cảm ơn Thầy Nguyễn Minh Trung đã hỗ trợ hướng dẫn em hoàn thành niên
luận.

Do trình độ lý luận và kinh nghiệm thực tế còn hạn chế nên bài báo cáo còn
nhiều thiếu sót, em rất mong nhận được ý kiến đóng góp từ Thầy, Cô để em có thể
rút ra nhiều kinh nghiệm hơn.

Cuối cùng, em xin chân thành cảm ơn gia đình và bạn bè, đã luôn tạo điều kiện,
quan tâm, giúp đỡ, động viên em trong suốt quá trình học tập và hoàn thành bài
niên luận này.

\begin{flushright}
  Em xin chân thành cảm ơn
\end{flushright}

\ctutableofcontents
\ctulistoffigures
\ctulistoftables

\ctuheading{DANH MỤC TỪ VIẾT TẮT}
\begin{longtable}{|>{\raggedright\arraybackslash}p{0.24\textwidth}|>{\raggedright\arraybackslash}p{0.68\textwidth}|}
  \hline
  \ctuTableHead{Kí hiệu \& Chữ viết tắt} & \ctuTableHead{Diễn giải} \\
  \hline
  \endhead
  API & Application Programming Interface (Giao diện lập trình ứng dụng) \\
  \hline
  AWS & Amazon Web Services (Nền tảng điện toán đám mây của Amazon) \\
  \hline
  CD & Continuous Delivery / Continuous Deployment (Phân phối liên tục / Triển khai liên tục) \\
  \hline
  CI & Continuous Integration (Tích hợp liên tục) \\
  \hline
  CVEs & Common Vulnerabilities and Exposures (Cơ sở dữ liệu lỗ hổng bảo mật toàn cầu) \\
  \hline
  GHCR & GitHub Container Registry (Kho lưu trữ Container của GitHub) \\
  \hline
  JDK & Java Development Kit (Bộ công cụ phát triển phần mềm Java) \\
  \hline
  JSON & JavaScript Object Notation (Định dạng lưu trữ và trao đổi dữ liệu) \\
  \hline
  LTS & Long-Term Support (Phiên bản hỗ trợ dài hạn) \\
  \hline
  MVC & Model-View-Controller (Mô hình kiến trúc phần mềm) \\
  \hline
  OCI & Open Container Initiative (Tiêu chuẩn mở cho định dạng container) \\
  \hline
  QA & Quality Assurance (Đội ngũ/Kiểm thử chất lượng) \\
  \hline
  QEMU & Quick EMUlator (Trình giả lập) \\
  \hline
  RAM & Random Access Memory (Bộ nhớ truy cập ngẫu nhiên) \\
  \hline
  SQL & Structured Query Language (Ngôn ngữ truy vấn dữ liệu -- đề cập trong SQL Injection) \\
  \hline
  UX & User Experience (Trải nghiệm người dùng) \\
  \hline
  vCPU & Virtual Central Processing Unit (Nhân CPU ảo) \\
  \hline
  XSS & Cross-Site Scripting (Lỗ hổng bảo mật tập lệnh chéo trang) \\
  \hline
  YAML & YAML Ain't Markup Language (Ngôn ngữ cấu hình dữ liệu) \\
  \hline
\end{longtable}
\ctuplaintablecaption{Danh mục từ viết tắt}

% -----------------------------------------------------------------------------
% MAIN MATTER
% -----------------------------------------------------------------------------
\ctumainmatter

\ctuheading{TÓM TẮT}
Trong bối cảnh công nghiệp phần mềm hiện đại, phương pháp phát triển linh hoạt
(Agile) đang được áp dụng rộng rãi nhằm tăng tốc độ phát triển và chuyển giao
tính năng mới. Tuy nhiên, tần suất cập nhật liên tục này vô hình trung tạo ra
nút thắt lớn tại giai đoạn triển khai, gây áp lực lên đội ngũ vận hành và làm
tăng rủi ro sai sót hệ thống do các thao tác thủ công. Để giải quyết vấn đề
này, đề tài ``Xây dựng hệ thống CI/CD cho ứng dụng Web'' được thực hiện nhằm tối
ưu hóa vòng đời phát triển phần mềm thông qua việc áp dụng tự động hóa theo mô
hình DevOps.

Niên luận tập trung vào việc thiết lập một luồng Tích hợp và Phân phối/triển
khai liên tục (CI/CD pipeline) cho một ứng dụng Web phát triển bằng ngôn ngữ
Java. Hệ thống đề xuất sử dụng nền tảng GitHub Actions làm cốt lõi để tự động
hóa toàn bộ các tác vụ: biên dịch mã nguồn (Build), thực thi kịch bản kiểm thử
(Test), và đánh giá chất lượng bảo mật (Security Evaluation). Sau khi vượt qua
các bước kiểm định nghiêm ngặt, hệ thống sẽ tự động đóng gói ứng dụng vào các
container (Docker Build) và phân phối an toàn lên kho lưu trữ GitHub Container
Registry (GHCR). Kết quả của đề tài mang lại một giải pháp chuẩn hóa, giúp giảm
thiểu tối đa sự can thiệp của con người, rút ngắn thời gian phát hành và nâng
cao tính ổn định, an toàn cho hệ thống.

\ctuheading{GIỚI THIỆU}
\introsection{Đặt vấn đề}
Trong bối cảnh ngành công nghiệp phần mềm hiện đại, phương pháp phát triển phần
mềm Agile đang trở thành tiêu chuẩn cốt lõi tại các công ty công nghệ lớn nhằm
đáp ứng nhu cầu thay đổi nhanh chóng của thị trường. Thông qua việc chia nhỏ quá
trình phát triển thành các chu kỳ ngắn (Sprints), đội ngũ lập trình
(Developer/Dev) có thể liên tục chuyển giao các tính năng phần mềm mới một cách
nhanh chóng.

Tuy nhiên, tốc độ lập trình và tần suất chuyển giao gia tăng này tạo ra một nút
thắt lớn ở giai đoạn triển khai (Deployment). Nó đặt áp lực nặng nề lên đội ngũ
vận hành (Operations/Ops). Ngoài việc phải đáp ứng khối lượng lớn các yêu cầu
cập nhật liên tục từ đội Dev, team Ops còn phải chịu trách nhiệm đảm bảo thời
gian hoạt động (uptime), tính sẵn sàng và sự ổn định của hệ thống máy chủ tổng
thể.

Sự mâu thuẫn giữa tốc độ phát triển và sự ổn định trong vận hành bộc lộ rõ qua
những điểm yếu của quy trình làm việc thủ công truyền thống:
\begin{itemize}
  \item \textbf{Ở giai đoạn Phát Triển (Development):} Thông thường, các nhóm
  phát triển làm việc độc lập trong nhiều tuần với khối lượng mã nguồn (source
  code) đồ sộ. Khi toàn bộ mã nguồn này được gộp chung và đưa lên máy chủ để
  biên dịch (build) cùng một lúc, hệ thống thường xuyên phát sinh hàng loạt lỗi
  phần mềm và xung đột cấu hình. Việc phát hiện lỗi quá muộn ở giai đoạn cuối
  khiến các kỹ sư mất rất nhiều thời gian để rà soát và khắc phục sự cố.
  \item \textbf{Ở giai đoạn Triển khai (Deployment):} Việc đưa ứng dụng lên môi
  trường thực tế (Production) là một chuỗi các công việc phức tạp, bao gồm cung
  cấp máy chủ, cấu hình các thư viện phụ thuộc (dependencies), thay đổi thiết
  lập mạng và triển khai các gói phần mềm đóng gói (artifact như định dạng
  jar/war, zip/tar). Khi nhóm Ops liên tục nhận các yêu cầu triển khai thủ công
  nhưng lại thiếu đi tài liệu hướng dẫn cấu hình rõ ràng, rủi ro về sai sót của
  con người là rất lớn, dẫn đến nguy cơ gián đoạn dịch vụ và giảm trải nghiệm
  người dùng.
\end{itemize}

Để giải quyết triệt để sự mâu thuẫn này, văn hóa DevOps ra đời nhằm phá bỏ rào
cản truyền thống, giúp đội ngũ Development và Operations làm việc gắn kết chặt
chẽ với nhau. Trọng tâm của DevOps là áp dụng các công cụ tự động hóa
(Automation Tools) nhằm giảm thiểu thao tác thủ công và loại bỏ sai sót từ con
người. Cốt lõi của mô hình này là quy trình Tích hợp liên tục và Phân
phối/triển khai liên tục (CI/CD):
\begin{itemize}
  \item Tích hợp liên tục (Continuous Integration -- CI): Yêu cầu mọi thay đổi
  (commit/push) mã nguồn của lập trình viên phải đi kèm với quá trình xây dựng
  và kiểm thử tự động ngay lập tức. Việc này giúp phát hiện và cô lập lỗi (sự
  cố nền, xung đột) ngay từ giai đoạn đầu để chúng không nhân lên diện rộng. Sau
  khi vượt qua kiểm định, mã nguồn được đóng gói thành phần mềm hoàn chỉnh
  (artifact). CI hoạt động như một hệ thống kiểm tra tự động. Mỗi khi code được
  thay đổi và commit lên hệ thống, các bài kiểm tra tự động sẽ được chạy để
  kiểm tra xem code có hoạt động đúng và không gây lỗi hay không.
  \item Phân phối liên tục (Continuous Delivery): là quá trình tiếp nối CI,
  trong đó mọi phiên bản phần mềm đã vượt qua các bước build, kiểm thử và đánh
  giá bảo mật sẽ được tự động đóng gói, lưu trữ tại kho artifact hoặc container
  registry, đồng thời sẵn sàng để triển khai lên các môi trường như Testing,
  Staging hoặc Production. Với Continuous Delivery, việc triển khai lên
  Production có thể cần sự phê duyệt thủ công của người quản trị.
  \item Triển khai liên tục (Continuous Deployment): là mức độ tự động hóa cao
  hơn của Continuous Delivery. Sau khi mã nguồn vượt qua toàn bộ các bước kiểm
  định trong pipeline, phiên bản mới sẽ được tự động triển khai lên môi trường
  chạy thực tế mà không cần thao tác thủ công. Để bảo đảm an toàn, quy trình
  này cần kết hợp cơ chế kiểm tra tình trạng hoạt động của ứng dụng (health
  check), theo dõi nhật ký hệ thống, kiểm soát phiên bản image và khả năng quay
  lui về phiên bản ổn định trước đó khi xảy ra sự cố.
\end{itemize}

Xuất phát từ những bất cập thực tế và tiềm năng tối ưu hóa trên, đề tài ``Xây
dựng hệ thống CI/CD cho ứng dụng Web'' được thực hiện. Niên luận tập trung vào
việc tự động hóa toàn bộ vòng đời phát triển của một ứng dụng Web viết bằng
ngôn ngữ Java.

Hệ thống đề xuất sẽ sử dụng nền tảng GitHub Actions để thiết lập một luồng
CI/CD (pipeline) tự động hóa. Các tác vụ cốt lõi bao gồm: tự động biên dịch mã
nguồn (Build), chạy các kịch bản kiểm thử (Test), đánh giá chất lượng và bảo
mật (Security Evaluation), tự động đóng gói ứng dụng vào các container Docker
(Docker Build) và phân phối an toàn lên kho lưu trữ GitHub Container Registry
(GHCR). Việc triển khai thành công hệ thống này không chỉ giúp tối ưu hóa thời
gian phát hành phần mềm mà còn trang bị môi trường hiện đại, bám sát các tiêu
chuẩn thực tế đang được áp dụng tại doanh nghiệp.

Sau khi Docker image được xuất bản lên GitHub Container Registry, máy chủ Ubuntu
sẽ tự động nhận biết phiên bản image mới thông qua dịch vụ Watchtower.
Watchtower định kỳ kiểm tra image mới trên GHCR, tải image về máy chủ, dừng
container phiên bản cũ và khởi tạo container phiên bản mới theo cấu hình đã
định nghĩa. Hoàn thiện bước Continuous Deployment cho môi trường thực nghiệm.

\introsection{Mục tiêu đề tài}
Mục tiêu chính của đề tài là xây dựng thành công hệ thống tự động hóa quy trình
Tích hợp liên tục, Phân phối liên tục và Triển khai liên tục (CI/CD) cho ứng
dụng Web phát triển bằng ngôn ngữ Java. Hệ thống tự động hóa từ khâu biên dịch,
kiểm thử, đánh giá chất lượng và bảo mật mã nguồn, đóng gói Docker image, xuất
bản image lên GitHub Container Registry đến tự động cập nhật ứng dụng trên máy
chủ Ubuntu. Việc triển khai nhằm giảm thiểu sự can thiệp thủ công, rút ngắn
thời gian phát hành, tăng tính nhất quán giữa các môi trường và bảo đảm khả
năng vận hành ổn định của hệ thống.

\introsection{Đối tượng và phạm vi nghiên cứu}
\introsubsection{Đối tượng nghiên cứu}
Đề tài tập trung nghiên cứu các đối tượng cốt lõi sau:

Quy trình Tích hợp liên tục (Continuous Integration -- CI) và Phân phối liên tục
(Continuous Delivery/Deployment -- CD) trong quy trình phát triển phần mềm hiện
đại.

Các công nghệ và nền tảng hỗ trợ tự động hóa DevOps, cụ thể bao gồm: nền tảng
chạy luồng công việc tự động GitHub Actions, công nghệ ảo hóa cấp hệ điều hành
(Containerization) bằng Docker, và kho lưu trữ container GitHub Container
Registry (GHCR).

Mã nguồn và cấu trúc của một ứng dụng Web được phát triển dựa trên ngôn ngữ
Java.

\introsubsection{Phạm vi nghiên cứu}
Để đảm bảo tính khả thi và tập trung đúng vào khía cạnh vận hành (Ops), đề tài
được giới hạn trong các phạm vi sau:

Về mặt ứng dụng: Đề tài không đi sâu vào việc phát triển các nghiệp vụ
(business logic) phức tạp cho ứng dụng Web. Thay vào đó, đề tài sẽ sử dụng một
ứng dụng Web Java (sample app) làm mẫu để tập trung trong việc thiết kế, kiểm
thử và tối ưu hóa hệ thống CI/CD.

Về mặt quy trình tự động hóa: Luồng CI/CD được thiết lập trên nền tảng GitHub
Actions với các tác vụ: tự động biên dịch mã nguồn (Build), chạy kiểm thử
(Test), kiểm tra coding convention, đánh giá chất lượng và bảo mật (Security
Evaluation), đóng gói ứng dụng thành Docker image (Docker Build), và xuất bản
image lên GitHub Container Registry (GHCR). Sau khi image được xuất bản thành
công, dịch vụ Watchtower trên máy ảo Ubuntu sẽ tự động kiểm tra, tải phiên bản
image mới và khởi động lại container ứng dụng. Qua đó, hệ thống thực hiện
Continuous Deployment trong phạm vi môi trường thực nghiệm.

Về môi trường thực nghiệm: Hệ thống được triển khai và cấu hình trên nền tảng
hệ sinh thái của GitHub (bao gồm Source code repository, GitHub Actions runner,
và tính năng GHCR) kết hợp với Docker Engine trên môi trường giả lập. Ngoài ra
hệ thống còn có phần cấu hình Vagrant để tự động triển khai máy ảo Ubuntu để
chạy hệ thống web cho phần Continuous Deployment.

Đề tài chưa đi sâu vào các bài toán quản lý cụm máy chủ lớn (như Kubernetes)
hay thiết lập hạ tầng mạng phức tạp trên các nền tảng Cloud (AWS, Azure và
Google Cloud).

\introsection{Công nghệ sử dụng}
\introsubsection{Các công nghệ được sử dụng trong đề tài}
\begin{center}
  \begin{tabular}{|>{\raggedright\arraybackslash}p{0.40\textwidth}|>{\raggedright\arraybackslash}p{0.50\textwidth}|}
    \hline
    \ctuTableHead{Nhóm công nghệ} & \ctuTableHead{Công nghệ} \\
    \hline
    Quản lý mã nguồn & Git, GitHub \\
    \hline
    Công cụ CI/CD & GitHub Actions \\
    \hline
    Build và kiểm thử & Maven \\
    \hline
    Kiểm tra coding convention & Checkstyle \\
    \hline
    Phân tích chất lượng mã nguồn & Aqua Security \\
    \hline
    Container & Docker (bao gồm Watchtower) \\
    \hline
    Container registry & GitHub Container Registry (GHCR) \\
    \hline
    Môi trường triển khai & Ubuntu 24.04 trên nền Vagrant và VMware \\
    \hline
    Công cụ hỗ trợ & Docker CLI \\
    \hline
  \end{tabular}
\end{center}
\ctuplaintablecaption{Các công nghệ được sử dụng}

\introsubsection{Giải thích các từ chuyên môn}
\textbf{Continuous Delivery:} là quy trình tự động đưa mọi thay đổi mã nguồn đã
được build, kiểm thử và kiểm tra bảo mật thành một phiên bản có thể triển khai
được. Phiên bản này thường được lưu trữ trên kho artifact hoặc container
registry và sẵn sàng triển khai lên môi trường Testing, Staging hoặc
Production. Việc phát hành lên Production có thể yêu cầu phê duyệt thủ công.

\textbf{Continuous Deployment:} là quy trình tự động triển khai phiên bản phần
mềm đã vượt qua toàn bộ kiểm định lên môi trường chạy thực tế mà không cần can
thiệp thủ công. Trong đề tài này, Continuous Deployment được thực hiện bằng
Watchtower: dịch vụ này theo dõi Docker image mới trên GHCR, tự động tải image
mới, thay thế container đang chạy và khởi động phiên bản ứng dụng mới trên máy
ảo Ubuntu.

\textbf{Watchtower:} là công cụ hỗ trợ tự động cập nhật các Docker container.
Watchtower định kỳ kiểm tra Docker registry để phát hiện image mới; khi có
phiên bản mới, công cụ sẽ tải image, dừng và xóa container cũ, sau đó khởi tạo
lại container mới dựa trên cấu hình hiện có. Công cụ này phù hợp với môi trường
thực nghiệm có quy mô nhỏ hoặc một máy chủ đơn lẻ.

\textbf{Vagrant:} là công cụ mã nguồn mở hỗ trợ tự động hóa việc tạo, cấu hình
và quản lý môi trường máy ảo. Vagrant sử dụng tệp cấu hình có tên Vagrantfile
để mô tả cấu hình máy ảo như hệ điều hành, dung lượng RAM, số lượng CPU, thiết
lập mạng và các bước cài đặt phần mềm. Trong đề tài, Vagrant được sử dụng để tự
động khởi tạo máy ảo Ubuntu làm môi trường triển khai ứng dụng Web, cài đặt
Docker Engine, Docker Compose và Watchtower. Nhờ đó, môi trường triển khai có
thể được tái tạo nhất quán, nhanh chóng và giảm các sai sót khi cấu hình thủ
công.

\textbf{GitHub Container Registry và Docker Hub:} hai nền tảng lưu trữ, quản lý
và phân phối container image (thường là Docker image) phổ biến nhất hiện nay.
Về cơ bản, chúng đóng vai trò như các ``kho'' chứa mã ứng dụng đã được đóng gói
cùng toàn bộ môi trường chạy, giúp việc triển khai phần mềm diễn ra đồng nhất
và nhanh chóng.

\textbf{Docker Hub:} Đây là kho lưu trữ container chính thức và lâu đời nhất
của Docker Inc. Nó được cộng đồng phần mềm sử dụng rộng rãi làm điểm đến mặc
định để tìm kiếm các ứng dụng và cơ sở hạ tầng dựng sẵn (như Nginx, MySQL,
Redis).

\textbf{GitHub Container Registry (GHCR):} Đây là dịch vụ lưu trữ container
tích hợp sẵn trong hệ sinh thái GitHub, cho phép nhà phát triển lưu trữ image
ngay tại nơi họ quản lý mã nguồn. Ưu điểm nổi bật của GHCR là kết nối cực kỳ
liền mạch với GitHub Actions để tự động hóa quy trình CI/CD, phân quyền quản lý
chi tiết (granular permissions).

\textbf{GitHub Actions:} nền tảng tự động hóa tích hợp sẵn trên \textbf{GitHub}.
Nó cho phép tạo các quy trình làm việc (workflow) để tự động hóa mọi bước trong
quy trình phát triển phần mềm, bao gồm build, test và triển khai ứng dụng ngay
khi có sự kiện (như push code hoặc tạo pull request) xảy ra.

\introsection{Nội dung nghiên cứu}
Để hoàn thành các mục tiêu đề ra, đề tài được tiến hành theo trình tự các bước
sau, tập trung sâu vào việc thiết lập, cấu hình và tối ưu hóa hệ thống vận
hành:

\introsubsection{Nghiên cứu cơ sở lý thuyết}
Tìm hiểu tổng quan về phương pháp Agile và văn hóa DevOps trong phát triển phần
mềm.

Nghiên cứu nguyên lý hoạt động, ưu nhược điểm của quy trình Tích hợp và Phân
phối/triển khai liên tục (CI/CD).

Khảo sát kiến trúc kỹ thuật của công nghệ Docker và nền tảng GitHub Actions.

\introsubsection{Phân tích và thiết kế quy trình (Pipeline Architecture)}
Thiết kế sơ đồ luồng công việc tự động (workflow) cho ứng dụng Java.

Xác định thứ tự các giai đoạn cốt lõi trong pipeline gồm: Checkout mã nguồn,
Thiết lập môi trường (Setup JDK), Biên dịch (Build), Kiểm thử (Test), Kiểm tra
coding convention, Quét bảo mật (Security Evaluation), Đóng gói Docker image
(Docker Build), Xuất bản image lên GHCR (Publish to GHCR), và tự động triển khai
phiên bản image mới lên máy ảo Ubuntu (Continuous Deployment) thông qua
Watchtower.

\introsubsection{Triển khai và cấu hình hệ thống}
Khởi tạo môi trường kho lưu trữ (repository) trên hệ sinh thái GitHub và đưa mã
nguồn ứng dụng mẫu lên hệ thống.

Xây dựng kịch bản Dockerfile đa bước (multi-stage) để tối ưu hóa dung lượng ứng
dụng Java.

Viết mã kịch bản cấu hình YAML (\texttt{.github/workflows/}) để tự động hóa các
tác vụ.

Tích hợp công cụ Security Evaluation để tự động phân tích mã nguồn và thiết lập
cơ chế chặn luồng (fail pipeline) nếu phát hiện lỗ hổng nghiêm trọng.

Thiết lập quản lý biến môi trường bí mật (GitHub Secrets) để phân quyền an toàn
khi hệ thống tự động đẩy Docker Image lên kho GHCR.

Tự động phát hiện thay đổi ở mã nguồn và tự động cập nhật bằng dịch vụ
Watchtower của Docker.

Cấu hình Docker Compose hoặc Docker CLI trên máy ảo Ubuntu để khởi tạo
container ứng dụng và dịch vụ Watchtower.

Cấu hình Watchtower xác thực với GHCR, định kỳ kiểm tra Docker image mới và tự
động cập nhật container ứng dụng khi phát hiện phiên bản mới được xuất bản.

Thiết lập health check cho container ứng dụng nhằm kiểm tra trạng thái hoạt
động sau khi triển khai; đồng thời lưu trữ tag phiên bản Docker image để hỗ trợ
xác định và quay lại phiên bản ổn định khi cần thiết.

\introsubsection{Vận hành thử nghiệm và Đánh giá}
Giả lập các thao tác của lập trình viên (commit, push code, tạo pull request)
để kích hoạt tự động luồng CI/CD.

Giám sát và kiểm tra tính đúng đắn của các log kiểm thử, báo cáo quét lỗi bảo
mật.

Xác minh tính khả dụng của Docker Image đầu ra trên GitHub Container Registry.

Tự động hoá việc cài đặt máy ảo Ubuntu để khởi tạo webserver và tự động phát
hiện ra thay đổi ở mã nguồn và tự động cập nhật mã nguồn.

Tổng hợp kết quả, đánh giá ưu/nhược điểm của hệ thống đề xuất so với phương
pháp thủ công và đưa ra hướng phát triển mở rộng trong tương lai.

\introsection{Bố cục của niên luận}
Niên luận được chia làm 4 chương:
\begin{itemize}
  \item Chương I: Đặc tả yêu cầu
  \item Chương II: Cơ sở lý thuyết
  \item Chương III: Thiết kế và cài đặt hệ thống
  \item Chương IV: Đánh giá kiểm thử
\end{itemize}

% -----------------------------------------------------------------------------
% CHƯƠNG I
% -----------------------------------------------------------------------------
\chapter{ĐẶC TẢ YÊU CẦU}
\label{chap:requirements}
\section{Tổng quan về giải pháp}
Để đáp ứng được tốc độ phát triển phần mềm hiện đại và sự ổn định trong vận
hành, giải pháp đề xuất tập trung vào việc tự động hóa toàn bộ vòng đời phát
triển của một ứng dụng Web viết bằng ngôn ngữ Java. Cốt lõi của giải pháp là áp
dụng văn hóa DevOps cùng các công cụ tự động hóa nhằm giảm thiểu thao tác thủ
công.

Cụ thể, hệ thống đề xuất sẽ sử dụng nền tảng GitHub Actions để thiết lập một
luồng Tích hợp liên tục và Phân phối liên tục (CI/CD pipeline) tự động hóa. Quy
trình giải pháp được thiết kế với các giai đoạn cốt lõi như sau:
\begin{itemize}
  \item Tự động hóa hoàn toàn các bước từ biên dịch mã nguồn (Build), chạy các
  kịch bản kiểm thử (Test) cho đến đánh giá chất lượng và bảo mật (Security
  Evaluation).
  \item Sau khi mã nguồn vượt qua các khâu kiểm định, luồng công việc sẽ tự động
  đóng gói ứng dụng vào các container Docker (Docker Build).
  \item Cuối cùng, các phiên bản đóng gói (Docker image) sẽ được phân phối và lưu
  trữ an toàn lên kho lưu trữ GitHub Container Registry (GHCR).
  \item Sau khi Docker image được xuất bản lên GHCR, dịch vụ Watchtower chạy trên
  máy ảo Ubuntu sẽ tự động phát hiện image mới, tải image về và thay thế
  container ứng dụng đang hoạt động. Bước này hiện thực hóa Continuous
  Deployment, giúp phiên bản mới được đưa vào môi trường chạy mà không yêu cầu
  quản trị viên thực hiện thủ công từng lần triển khai.
\end{itemize}

Thông qua giải pháp này, đề tài xây dựng một hệ thống chuẩn hóa, giúp loại bỏ
các rào cản truyền thống giữa nhóm phát triển (Dev) và nhóm vận hành (Ops), đồng
thời mang lại một môi trường triển khai đồng nhất, độc lập và bám sát các tiêu
chuẩn công nghiệp.

\section{Các vấn đề cần giải quyết}
Hệ thống CI/CD được thiết kế nhằm giải quyết trực diện những điểm yếu của quy
trình làm việc thủ công truyền thống. Dưới đây là các vấn đề trọng tâm mà hệ
thống cần khắc phục:
\begin{itemize}
  \item \textbf{Sự chậm trễ và xung đột trong tích hợp mã nguồn:} Thông thường,
  các nhóm phát triển làm việc độc lập trong nhiều tuần với khối lượng mã nguồn
  đồ sộ. Việc gộp chung mã nguồn và đưa lên máy chủ để biên dịch cùng lúc
  thường xuyên làm phát sinh hàng loạt lỗi phần mềm và xung đột cấu hình. Vấn
  đề này khiến kỹ sư mất rất nhiều thời gian để rà soát và khắc phục sự cố do
  phát hiện lỗi quá muộn.
  \item \textbf{Rủi ro từ các thao tác triển khai thủ công:} Quá trình đưa ứng
  dụng lên môi trường thực tế (Production) là một chuỗi các công việc phức tạp
  như cung cấp máy chủ, cấu hình thư viện phụ thuộc, và triển khai các gói
  artifact (jar/war, zip/tar). Khi các thao tác này được thực hiện thủ công mà
  thiếu tài liệu cấu hình rõ ràng, rủi ro sai sót của con người là rất lớn, dễ
  dẫn đến gián đoạn dịch vụ và giảm trải nghiệm người dùng.
  \item \textbf{Áp lực quá tải lên đội ngũ vận hành (Ops):} Tốc độ phát triển
  nhanh chóng (thông qua các Sprints) tạo ra nút thắt lớn ở giai đoạn triển
  khai. Điều này đặt áp lực nặng nề lên đội ngũ Operations, khi họ vừa phải đáp
  ứng khối lượng lớn yêu cầu cập nhật liên tục từ nhóm Dev, vừa phải chịu trách
  nhiệm đảm bảo thời gian hoạt động (uptime) và sự ổn định của hệ thống máy chủ
  tổng thể.
\end{itemize}

% -----------------------------------------------------------------------------
% CHƯƠNG II
% -----------------------------------------------------------------------------
\chapter{CƠ SỞ LÝ THUYẾT}
\label{chap:theory}
\section{Tổng quan về CI/CD}
CI/CD là một trong những tiến bộ quan trọng trong phát triển phần mềm hiện đại,
giúp thay đổi cách các nhóm DevOps quản lý quá trình phát triển và triển khai
sản phẩm. Ban đầu, các quy trình phát triển phần mềm thường bị chia tách với
quy trình kiểm thử và triển khai, dẫn đến phát hiện lỗi muộn và khó khăn trong
việc duy trì tính ổn định khi tích hợp nhiều thay đổi cùng lúc.

CI/CD (Continuous Integration/Continuous Delivery -- Deployment) là một tập hợp
các phương pháp và công cụ giúp tự động hóa quy trình phát triển, kiểm thử, và
triển khai phần mềm. Mục tiêu chính của CI/CD là đẩy nhanh tốc độ phát hành
phần mềm, giảm lỗi và nâng cao tính ổn định của các phiên bản phát hành
(release).

\ctufigure{image2.png}{1.0}{Quy trình CI/CD}{fig:cicd-process}

\subsection{Tích hợp liên tục}
CI (Continuous Integration) là quá trình liên tục kiểm tra và tích hợp code mới
vào code chính của một dự án. Mục tiêu của CI là đảm bảo rằng bất kỳ thay đổi
nào trong code (chẳng hạn như sửa lỗi hoặc thêm tính năng) đều được kiểm tra
ngay lập tức để phát hiện lỗi trước khi nó gây ảnh hưởng xấu đến hệ thống.

CI hoạt động như một hệ thống kiểm tra tự động. Mỗi khi bạn thay đổi code và
commit lên hệ thống, các bài kiểm tra tự động sẽ được chạy để kiểm tra xem code
có hoạt động đúng và không gây lỗi hay không.

\subsubsection{Những bước cơ bản trong CI}
\begin{enumerate}
  \item \textbf{Functional tests:} Kiểm tra xem phần mềm có hoạt động đúng theo
  mong đợi hay không. Đây thường là các bài kiểm tra tự động mà đội ngũ QA
  (Quality Assurance) phát triển để đảm bảo tính năng hoạt động đúng. VD: Người
  dùng có thể thêm sản phẩm vào giỏ hàng và thực hiện thanh toán thành công hay
  không.
  \item \textbf{Security scans:} Kiểm tra code có lỗ hổng bảo mật nào không,
  chẳng hạn như khả năng bị SQL Injection hay
  \href{https://200lab.io/blog/giai-phap-bao-mat-giao-dien-web-ngan-chan-tan-cong-xss-va-csrf-hieu-qua/}{XSS}
  (Cross-Site Scripting).
  \item \textbf{Code quality scans:} Đảm bảo code tuân thủ các tiêu chuẩn, ví dụ
  như độ dài hàm, cách sử dụng khoảng trắng, và các quy tắc coding style. VD:
  Tiêu chuẩn như đặt tên biến.
  \item \textbf{Performance tests:} Kiểm tra xem code có đáp ứng được yêu cầu về
  hiệu suất không, chẳng hạn như thời gian xử lý yêu cầu hoặc khả năng chịu
  tải.
  \item \textbf{License scanning:} Kiểm tra xem tất cả các thư viện hoặc công cụ
  bạn sử dụng có giấy phép phù hợp hay không, để tránh các vấn đề pháp lý.
  \item \textbf{Fuzz testing:} Gửi các dữ liệu bất thường, như chuỗi ký tự quá
  dài hoặc số không hợp lệ, vào ứng dụng để kiểm tra xem nó có bị crash hay gặp
  lỗi bất ngờ không.
\end{enumerate}

\subsubsection{Lợi ích của việc thiết lập CI}
CI giúp phát hiện sớm các lỗi tiềm ẩn, tiết kiệm rất nhiều thời gian và công
sức, rõ ràng việc sửa một lỗi ngay khi vừa phát hiện sẽ dễ dàng hơn nhiều so
với khi lỗi đó đã tích hợp vào nhiều phần khác của hệ thống.

Ví dụ: Giả sử đội ngũ lập trình đang phát triển một tính năng mới cho ứng dụng.
Thay vì đợi đến cuối quá trình phát triển để kiểm tra, CI cho phép bạn thực
hiện các bài kiểm tra ngay khi lập trình viên viết code. Nhờ đó, nếu có lỗi
xuất hiện, họ có thể phát hiện ngay lập tức và sửa chữa khi vấn đề còn nhỏ,
thay vì đợi đến khi tính năng đã hoàn thành và tích hợp với các phần khác.

CI giúp toàn bộ đội nhóm đều có thể thấy tình trạng hiện tại của phần mềm, tất
cả các thành viên từ lập trình viên, quản lý, đến QA (kiểm thử chất lượng), UX
(trải nghiệm người dùng), và thậm chí cả bộ phận bảo mật có thể theo dõi tiến
trình và trạng thái của dự án. Từ đó, mọi người có thể điều chỉnh công việc của
mình để phù hợp với tình trạng thực tế.

Ví dụ: Khi các bài kiểm tra bảo mật được thực hiện thường xuyên, cả nhóm sẽ
nhận biết được dự án có gặp phải các vấn đề bảo mật mới không. Nếu CI phát hiện
rằng một đoạn code mới làm lộ lỗ hổng bảo mật, thì không chỉ các lập trình viên
mà cả các quản lý dự án cũng biết ngay lập tức. Nhóm sẽ có thể phân bổ nguồn
lực để sửa lỗi hoặc điều chỉnh kế hoạch phát triển.

\subsection{Phân phối/triển khai liên tục}
CD là quá trình đưa code từ nơi bạn viết lên một nơi mà người dùng có thể sử
dụng được, như website hoặc ứng dụng di động. Có 2 loại CD chính:
\begin{itemize}
  \item \textbf{Continuous Delivery (Phân phối liên tục):} Đảm bảo rằng code
  luôn sẵn sàng để được triển khai bất cứ lúc nào. Sau khi code đã vượt qua các
  kiểm tra CI, nó được đặt ở trạng thái sẵn sàng để triển khai. Tuy nhiên, việc
  triển khai có thể được thực hiện thủ công.
  \item \textbf{Continuous Deployment (Triển khai liên tục):} Tự động triển khai
  code ngay khi nó vượt qua tất cả các kiểm tra. Không cần sự can thiệp thủ
  công, hệ thống sẽ tự động đưa code vào môi trường production.
\end{itemize}

\subsubsection{Những bước cơ bản trong CD}
Build: Trước khi triển khai code, thực hiện quá trình build cho các ngôn ngữ
biên dịch (compiled language) như C, C++, Java,~\ldots

Deploy: Triển khai code: đẩy Docker image lên repository, dùng dòng lệnh để đưa
code lên AWS,~\ldots{} Triển khai có thể được tự động hoặc thủ công, tùy thuộc
vào cách bạn thiết lập.

Ví dụ: Khi hoàn thành thêm tính năng mới như giỏ hàng, CD sẽ tự động kiểm tra
và đóng gói tính năng này rồi triển khai lên website thật để khách hàng có thể
sử dụng ngay lập tức.

\subsubsection{Lợi ích của việc thiết lập CD}
Mục tiêu của CD là giúp việc triển khai và phát hành (release) phần mềm trở nên
dễ dàng và có thể diễn ra thường xuyên mà không gặp nhiều rủi ro. Khi CD được
thiết lập, mỗi lần bạn commit code, code có thể tự động được triển khai vào các
môi trường như: review environment, Staging environment hay production
environment.

Các môi trường khác ngoài production (như review và staging) hoạt động giống
như những phép thử, nhằm đảm bảo code chạy ổn định trước khi release chính thức
cho người dùng.

\section{Nền tảng tự động hóa GitHub Actions}
GitHub Actions là nền tảng tự động hóa được tích hợp sẵn trực tiếp trên hệ sinh
thái GitHub. Nền tảng này cho phép tạo các quy trình làm việc (workflow) nhằm
tự động hóa mọi bước trong vòng đời phát triển phần mềm, từ build, test cho đến
deploy ứng dụng ngay khi có sự kiện xảy ra (như push code hoặc tạo pull
request).

Trong kiến trúc của GitHub Actions, một luồng công việc tự động được định nghĩa
thông qua tệp tin cấu hình YAML và bao gồm các thành phần:
\begin{itemize}
  \item \textbf{Events (Sự kiện):} Các trigger kích hoạt quy trình, ví dụ như
  \texttt{push} hoặc \texttt{pull\_request} trên nhánh \texttt{main}.
  \item \textbf{Jobs (Công việc):} Một luồng có thể chạy một hoặc nhiều job (như
  Build, Testing, Security\_Scan, Publish) tuần tự hoặc song song trên các máy
  chủ ảo (runner) do GitHub cung cấp (ví dụ: \texttt{ubuntu-latest}).
  \item \textbf{Steps (Bước):} Từng nhiệm vụ cụ thể bên trong một job, có thể là
  thực thi một lệnh shell hoặc gọi các actions có sẵn (như
  \texttt{actions/checkout} để kéo mã nguồn).
  \item \textbf{Artifacts (Tài nguyên đầu ra):} Chức năng lưu trữ tạm thời các
  gói phần mềm hoặc tệp báo cáo (như file \texttt{.war} hoặc báo cáo bảo mật
  JSON) để chia sẻ giữa các job với nhau.
\end{itemize}

\ctufigure{image3.png}{0.58}{Kiến trúc của GitHub Actions}{fig:gha-architecture}

\section{Công cụ quản lý dự án và kiểm tra mã nguồn Java}
Ứng dụng Web trong đề tài được phát triển dựa trên ngôn ngữ Java, do đó yêu cầu
các công cụ đặc thù để quản lý vòng đời và chất lượng:
\begin{itemize}
  \item Apache Maven: Đây là công cụ quản lý dự án (quản lý thư viện phụ thuộc,
  quản lý vòng đời biên dịch) phổ biến nhất cho Java. Trong hệ thống, Maven
  đóng vai trò tải các thư viện, chạy các kịch bản kiểm thử (\texttt{mvn test})
  và đóng gói mã nguồn thành phần mềm hoàn chỉnh định dạng WAR
  (\texttt{mvn install}).
  \item Checkstyle: Để đảm bảo tính đồng nhất trong phong cách lập trình và phát
  hiện sớm các đoạn code không tuân thủ tiêu chuẩn, công cụ Checkstyle được
  tích hợp vào bước kiểm thử (\texttt{mvn checkstyle:checkstyle}) làm rào chắn
  (coding convention validator).
\end{itemize}

\section{Đánh giá bảo mật (Security Evaluation)}
Đề tài sử dụng \textbf{Trivy (Aqua Security)}, một công cụ quét lỗ hổng bảo mật
mã nguồn mở toàn diện. Trivy phân tích cấu trúc tệp tin (filesystem scan) của
toàn bộ dự án để đối chiếu hệ điều hành và các thư viện phụ thuộc với các cơ sở
dữ liệu lỗ hổng bảo mật toàn cầu (CVEs). Kết quả quét được xuất ra dưới định
dạng JSON và có cấu hình chặn đứng toàn bộ tiến trình CI/CD nếu phát hiện lỗ
hổng nghiêm trọng.

\section{Công nghệ ảo hóa (Docker)}
Docker là nền tảng phần mềm cho phép các nhà phát triển xây dựng, kiểm thử và
triển khai các ứng dụng thông qua công nghệ ảo hóa ở dạng container. Docker
đóng gói ứng dụng Web Java cùng toàn bộ môi trường chạy thực tế (runtime, thư
viện hệ thống) thành các đơn vị chuẩn hóa gọi là Container. Nhờ vậy, ứng dụng
có thể chạy đồng nhất trên bất kỳ hạ tầng nào mà không gặp tình trạng xung đột
môi trường. Trong dự án, quá trình này được thực thi thông qua kịch bản
Dockerfile đa bước (multi-stage) và lệnh \texttt{docker build}.

\section{GitHub Container Registry (GHCR)}
GitHub Container Registry và Docker Hub là hai nền tảng lưu trữ, quản lý và
phân phối container image phổ biến. Tuy nhiên, thay vì sử dụng AWS ECR hay
Docker Hub, đề tài lựa chọn GHCR vì đây là dịch vụ được tích hợp sẵn trong hệ
sinh thái GitHub.

Việc sử dụng GHCR mang lại những lợi ích cốt lõi cho hệ thống:
\begin{itemize}
  \item Cho phép nhà phát triển lưu trữ Docker image ngay tại nơi họ đang quản lý
  mã nguồn.
  \item Kết nối cực kỳ liền mạch và có độ tin cậy cao với GitHub Actions.
  \item Đảm bảo an toàn bảo mật thông qua cơ chế phân quyền chi tiết (granular
  permissions) và việc xác thực tự động không cần chia sẻ mật khẩu tĩnh
  (secretless authentication) thông qua token mặc định của GitHub.
\end{itemize}

\section{Watchtower và cơ chế tự động cập nhật container}
Watchtower là công cụ hỗ trợ tự động cập nhật Docker container. Công cụ này
chạy như một container trên máy chủ triển khai và định kỳ kiểm tra Docker
registry để phát hiện image mới của các container được theo dõi.

Trong hệ thống của đề tài, Watchtower được cấu hình trên máy ảo Ubuntu có cài
Docker Engine. Sau khi GitHub Actions xuất bản Docker image mới lên GHCR,
Watchtower sẽ xác thực với GHCR, tải image mới về máy chủ, dừng container ứng
dụng cũ và khởi động lại container mới từ image vừa tải.

Cơ chế này hoàn thiện giai đoạn Continuous Deployment trong phạm vi nghiên cứu.
Luồng triển khai tổng quát của hệ thống được mô tả như sau:

\begin{center}
  \begin{minipage}{0.92\textwidth}
    \raggedright
    Push source code $\rightarrow$ GitHub Actions
    build/test/checkstyle/security scan $\rightarrow$ Docker Build
    $\rightarrow$ Push Docker image lên GHCR $\rightarrow$ Watchtower phát hiện
    image mới $\rightarrow$ Pull image $\rightarrow$ Khởi động lại container
    $\rightarrow$ Health check ứng dụng.
  \end{minipage}
\end{center}

% -----------------------------------------------------------------------------
% CHƯƠNG III
% -----------------------------------------------------------------------------
\chapter{THIẾT KẾ VÀ CÀI ĐẶT GIẢI PHÁP}
\label{chap:design}
\section{Phân tích kiến trúc hệ thống ứng dụng Web mẫu}
Ứng dụng web mẫu trong đề tài (VProfile) được viết bằng ngôn ngữ lập trình Java
và được thiết kế theo mô hình đa tầng (Multi-tier) bao gồm các thành phần:
\begin{itemize}
  \item Tầng ứng dụng web (Web Tier): Sử dụng Nginx làm reverse proxy và cân bằng
  tải (load balancer).
  \item Tầng Ứng dụng (Application Tier): Được viết bằng ngôn ngữ Java và
  Framework Spring, chạy trên máy chủ Apache Tomcat (port 8080). Đồng thời đóng
  gói phần mềm đầu ra của tầng này là một tệp định dạng \texttt{.war}.
  \item Tầng Cơ sở dữ liệu và Dịch vụ hỗ trợ: Bao gồm MySQL 8.0, Memcached
  (caching), và RabbitMQ (Message Broker).
  \item Hệ thống CI/CD trong đề tài sẽ tập trung xử lý tự động hóa cho Tầng Ứng
  dụng web (Java Spring MVC), đảm bảo mã nguồn được kiểm thử, đóng gói thành
  file \texttt{.war}, và cuối cùng xây dựng thành Docker image để sẵn sàng tích
  hợp với các tầng dịch vụ khác.
\end{itemize}

\ctufigure{image4.png}{1.0}{Kiến trúc Web App Java}{fig:webapp-architecture}

\section{Thiết kế luồng công việc tự động hóa (Pipeline Design)}
Để tối ưu hóa thời gian chạy và đảm bảo tính chặt chẽ, pipeline được chia thành
5 công việc (Jobs) chính, hoạt động kết hợp giữa tuần tự và song song:
\begin{enumerate}
  \item Job CI (Build, Test \& Package): Chạy đầu tiên để kiểm thử mã nguồn và
  đóng gói thành một artifact (\texttt{.war}).
  \item Job Security\_Scan: Hoạt động song song với Job CI nhằm đánh giá bảo mật
  mà không làm kéo dài thời gian chờ đợi.
  \item Job BUILD\_AMD64 \& BUILD\_ARM64: Chỉ được kích hoạt khi cả hai khâu
  kiểm thử và bảo mật đều thành công. Hai job này chạy song song trên các hạ
  tầng khác nhau để đóng gói Docker cho cả hai kiến trúc vi xử lý (x86\_64 và
  ARM64).
  \item Job PUBLISH\_TO\_GHCR: Hoạt động cuối cùng, gộp kết quả của hai tiến
  trình build trên thành một container đa kiến trúc (Multi-arch) và phân phối
  lên kho lưu trữ của GitHub (GHCR).
  \item Sau khi image được xuất bản thành công, dịch vụ Watchtower trên máy ảo
  Ubuntu sẽ tự động kiểm tra, tải phiên bản image mới và khởi động lại container
  ứng dụng. Qua đó, hệ thống thực hiện Continuous Deployment trong phạm vi môi
  trường thực nghiệm.
\end{enumerate}

\section{Cài đặt quy trình CI (Continuous Integration)}
Quá trình cài đặt giải pháp được thực hiện thông qua tệp tin cấu hình YAML
(\texttt{.github/workflows/main.yml}) đặt tại kho lưu trữ mã nguồn.

\subsection{Tạo access token trên GitHub}
Tạo \textbf{Personal access tokens} để đăng nhập vào GHCR (GitHub Container
Registry).

Các quyền cần thiết:
\begin{itemize}
  \item Quyền quản lý repo
  \item Quyền cập nhật workflow
  \item Quyền ghi và đọc packages
\end{itemize}

\ctufigure{image5.png}{1.0}{Tạo personal access token (classic)}{fig:pat}

\subsection{Cấu hình sự kiện kích hoạt (Triggers)}
Hệ thống được thiết lập cơ chế tự động kích hoạt (trigger) khi có sự kiện đẩy
mã nguồn (\texttt{push}) hoặc tạo yêu cầu gộp mã (\texttt{pull\_request}) vào
nhánh chính (\texttt{main}). Nền tảng được cấp quyền \texttt{contents: read} để
sao chép mã nguồn và \texttt{packages: write} để đẩy các container lên GitHub
Container Registry (GHCR).

\ctufigure{image6.png}{0.42}{Cấu hình sự kiện kích hoạt}{fig:triggers}

\subsection{Cài đặt giai đoạn Tích hợp liên tục (CI)}
Giai đoạn Tích hợp liên tục chịu trách nhiệm biên dịch và kiểm thử ứng dụng
Java.
\begin{itemize}
  \item Thiết lập môi trường: Khởi tạo máy chủ ảo \texttt{ubuntu-latest}, cài
  đặt JDK 17 (Temurin). Đặc biệt, cơ chế bộ nhớ đệm (\texttt{cache: 'maven'})
  được kích hoạt để lưu lại các thư viện phụ thuộc, giúp giảm thiểu thời gian
  tải lại các package trong các lần chạy sau.
  \item Biên dịch và Kiểm thử: Lệnh \texttt{mvn install} được thực thi để
  compile, test, và package ra tệp \texttt{.war}.
  \item Kiểm tra mã nguồn: Lệnh \texttt{mvn checkstyle:checkstyle} được gọi để
  đảm bảo mã nguồn tuân thủ các quy tắc lập trình (coding convention).
  \item Lưu trữ Artifact đầu ra: Gói phần mềm \texttt{vprofile-v2.war} được
  trích xuất và tải lên bộ nhớ tạm của GitHub Actions (thời hạn 7 ngày) để tiết
  kiệm tài nguyên.
\end{itemize}

\ctufigure{image7.png}{1.0}{Job CI}{fig:job-ci}

\subsection{Cài đặt giai đoạn Đánh giá bảo mật DevSecOps}
Tuân thủ nguyên tắc ``Shift-Left Security'', công việc Security\_Scan được cấu
hình chạy song song độc lập với tiến trình CI.

Hệ thống sử dụng action \texttt{aquasecurity/trivy-action} để phân tích tĩnh
toàn bộ hệ thống tập tin (filesystem), tìm kiếm các lỗ hổng của thư viện và hệ
điều hành dựa trên dữ liệu CVE. Kết quả báo cáo được xuất ra dưới định dạng
JSON làm cơ sở đánh giá chất lượng.

\ctufigure{image8.png}{1.0}{Job Security\_Scan}{fig:job-security}

\subsection{Cài đặt giai đoạn Đóng gói đa kiến trúc (Multi-Arch Docker Build)}
Thay vì tự biên dịch lại mã nguồn bên trong Docker (gây dư thừa tài nguyên), hệ
thống áp dụng nguyên tắc ``Build Once'' (Chỉ biên dịch một lần).

Các Job đóng gói sẽ tải tệp \texttt{.war} từ bước CI về. Dockerfile chỉ cần
thực hiện lệnh \texttt{COPY} tệp tin này vào máy chủ Tomcat.

Để hỗ trợ tính tương thích cho nhiều nền tảng, mã nguồn được build đồng thời
cho hai kiến trúc:
\begin{itemize}
  \item AMD64: Thực thi trên runner mặc định \texttt{ubuntu-latest}.
\end{itemize}

\ctufigure{image9.png}{1.0}{Job BUILD\_AMD64}{fig:job-amd64}

\begin{itemize}
  \item ARM64: Thực thi trực tiếp trên hạ tầng Native ARM
  (\nolinkurl{ubuntu-24.04-arm}). Việc không sử dụng giả lập QEMU (Quick EMUlator)
  giúp tốc độ build phần mềm tăng lên đáng kể.
\end{itemize}

\ctufigure{image10.png}{1.0}{Job BUILD\_ARM64}{fig:job-arm64}

\begin{itemize}
  \item Thay vì gán nhãn \texttt{latest} ngay lúc này, cả hai Job đều đẩy (push)
  image lên GHCR dưới dạng chuỗi mã băm (Digest) để tránh xung đột dữ liệu.
\end{itemize}

\section{Cài đặt quy trình CD (Continuous Delivery)}
Sau khi các container đơn lẻ hoàn tất, công việc PUBLISH\_TO\_GHCR sử dụng công
cụ \texttt{docker buildx imagetools} để ghép nối hai chuỗi mã băm (Digest của
AMD64 và ARM64) thành một OCI Manifest duy nhất.

Kết quả đầu ra là một Docker Image hoàn chỉnh trên kho lưu trữ GHCR. Người dùng
máy chủ truyền thống hay các dòng vi xử lý kiến trúc ARM (như Apple Silicon)
khi dùng \texttt{docker pull} ứng dụng đều sẽ nhận được container tương thích
tự động.

\ctufigure{image11.png}{1.0}{Job PUBLISH\_TO\_GHCR}{fig:job-publish}

\section{Cài đặt quy trình CD (Continuous Deployment)}
\subsection{Cài đặt máy ảo Ubuntu 24.04 bằng Vagrant}
\textbf{Vagrant được cấu hình như sau:}
\begin{itemize}
  \item \textbf{\texttt{Vagrant.configure("2")}:} Sử dụng cấu hình theo chuẩn
  phiên bản API v2 của Vagrant.
  \item \textbf{\texttt{config.vm.box = "bento/ubuntu-24.04"}:} Sử dụng hệ điều
  hành Ubuntu 24.04 LTS được đóng gói bởi Bento.
  \item \textbf{\texttt{virtualbox.memory = 4096}:} Cấp phát 4096 MB (4 GB) bộ
  nhớ RAM cho máy ảo để đáp ứng nhu cầu build image và chạy nhiều container
  cùng lúc.
  \item \textbf{\texttt{virtualbox.cpus = 4}:} Phân bổ 4 nhân CPU ảo (vCPU) giúp
  tăng tốc độ build và xử lý tác vụ của ứng dụng.
  \item \textbf{Cổng 80 (Guest) $\rightarrow$ 8080 (Host):} Cho phép truy cập
  vào Reverse Proxy (ví dụ: Nginx) từ trình duyệt máy thật thông qua địa chỉ
  \url{http://localhost:8080}.
  \item \textbf{Cổng 8080 (Guest) $\rightarrow$ 8081 (Host):} Cho phép truy cập
  trực tiếp vào dịch vụ backend/ứng dụng qua địa chỉ
  \url{http://localhost:8081}.
  \item \texttt{auto\_correct: true}: Tự động chuyển sang một cổng khả dụng khác
  trên máy Host nếu cổng được chỉ định (8080 hoặc 8081) đang bị ứng dụng khác
  chiếm dụng.
  \item \texttt{config.vm.provision}: tự động hoá cài đặt và khởi tạo dịch vụ
  Docker.
\end{itemize}

\ctufigure{image12.png}{1.0}{Cấu hình Vagrant}{fig:vagrant}

\begin{table}[!htbp]
  \centering
  \fontsize{11pt}{13.42pt}\selectfont
  \setlength{\tabcolsep}{3pt}
  \renewcommand{\arraystretch}{1.25}
  \begin{tabularx}{\textwidth}{|>{\hsize=0.8\hsize\linewidth=\hsize\raggedright\arraybackslash}X|>{\hsize=1.1\hsize\linewidth=\hsize\raggedright\arraybackslash}X|>{\hsize=1.1\hsize\linewidth=\hsize\raggedright\arraybackslash}X|}
    \hline
    \ctuTableHead{Thành phần} & \ctuTableHead{Cấu hình} & \ctuTableHead{Mục đích / Chức năng} \\
    \hline
    Virtualization Provider & VMware & Nền tảng ảo hóa máy ảo \\
    \hline
    Base OS & bento/ubuntu-24.04 (v202510.26.0) & Hệ điều hành Ubuntu 24.04 LTS \\
    \hline
    CPU / RAM & 4 vCPUs / 4096 MB RAM & Đảm bảo hiệu năng cho Docker và các dịch vụ \\
    \hline
    Port 80 (Guest) & 8080 (Host) & Truy cập Reverse Proxy (\url{http://localhost:8080}) \\
    \hline
    Port 8080 (Guest) & 8081 (Host) & Truy cập trực tiếp Backend (\url{http://localhost:8081}) \\
    \hline
    Shared Directory & /vagrant $\leftrightarrow$ Thư mục dự án & Đồng bộ mã nguồn theo thời gian thực \\
    \hline
    Container Engine & docker.io, docker-compose-v2 & Quản lý và thực thi môi trường container \\
    \hline
    Registry Authentication & GitHub Container Registry (ghcr.io) & Xác thực để kéo/đẩy container image \\
    \hline
  \end{tabularx}
  \caption{Bảng cấu hình Vagrant}
  \label{tab:vagrant}
\end{table}
\FloatBarrier

\subsection{Cài đặt docker-compose}
\subsubsection{Tổng quan kiến trúc hệ thống}
\textbf{Hệ thống được chia thành các lớp chức năng sau:}
\begin{itemize}
  \item \textbf{Lớp Web (Presentation/Proxy Layer):} vproweb (Nginx).
  \item \textbf{Lớp Ứng dụng (Application Layer):} vproapp (Java Web/Backend
  App).
  \item \textbf{Lớp Middleware \& Cache:} vpromq01 (RabbitMQ), vprocache01
  (Memcached).
  \item \textbf{Lớp Dữ liệu (Database Layer):} vprodb (MySQL Database).
  \item \textbf{Lớp Tự động hóa CD (Continuous Deployment Layer):} watchtower (Tự
  động cập nhật phiên bản mới).
\end{itemize}

\ctufigure{image13.png}{1.0}{Sơ đồ kiến trúc hệ thống Web trên nền Docker}{fig:docker-architecture}

\subsubsection{Cấu hình dịch vụ Cơ sở dữ liệu vprodb}
\textbf{Dịch vụ cơ sở dữ liệu vprodb được cấu hình như sau:}
\begin{itemize}
  \item \textbf{build.context:} Xây dựng image từ Dockerfile trong thư mục
  \nolinkurl{./Docker-files/db} (chứa các script khởi tạo schema database ban
  đầu).
  \item \textbf{ports (3306:3306):} Mở cổng chuẩn của MySQL để cho phép kết nối
  từ trong mạng container hoặc từ bên ngoài.
  \item \textbf{volumes (vprodbdata:/var/lib/mysql):} Gắn volume vprodbdata vào
  thư mục lưu trữ dữ liệu của MySQL. Đảm bảo tính toàn vẹn và bền vững của dữ
  liệu (Data Persistence), dữ liệu không bị mất khi container bị xóa hoặc khởi
  động lại.
  \item \textbf{environment:} Thiết lập mật khẩu lấy từ biến môi trường cho tài
  khoản root của cơ sở dữ liệu.
\end{itemize}

\ctufigure{image14.png}{0.83}{Cơ sở dữ liệu vprodb}{fig:vprodb}

\subsubsection{Cấu hình dịch vụ Bộ nhớ đệm vprocache01}
\textbf{Dịch vụ Bộ nhớ đệm vprocache01 được cấu hình như sau:}
\begin{itemize}
  \item \textbf{image: memcached:} Sử dụng image Memcached chính thức từ Docker
  Hub.
  \item \textbf{Mục đích:} Cung cấp cơ chế lưu trữ đệm trong RAM (In-memory
  Caching) tại cổng 11211, giúp giảm tải truy vấn cho cơ sở dữ liệu và tăng tốc
  độ phản hồi cho người dùng.
\end{itemize}

\ctufigure{image15.png}{0.40}{Bộ nhớ đệm vprocache01}{fig:vprocache}

\subsubsection{Cấu hình dịch vụ RabbitMQ}
\textbf{Dịch vụ RabbitMQ được cấu hình như sau:}
\begin{itemize}
  \item \textbf{image rabbitmq:} Sử dụng hệ thống Message Broker RabbitMQ.
  \item \textbf{ports (5672:5672):} Cổng giao thức AMQP dùng để truyền/nhận các
  bản tin (message) giữa các thành phần.
  \item \textbf{Mục đích:} Hỗ trợ xử lý bất đồng bộ (Asynchronous Processing) và
  phân tách phụ thuộc (Decoupling) giữa các module tác vụ.
\end{itemize}

\ctufigure{image16.png}{0.84}{Cấu hình dịch vụ RabbitMQ}{fig:rabbitmq}

\subsubsection{Cấu hình dịch vụ Watchtower}
Dịch vụ Watchtower được cấu hình nhằm tự động phát hiện và triển khai image mới
của ứng dụng Backend.
\begin{itemize}
  \item \textbf{image:} Sử dụng image \texttt{nickfedor/watchtower:1.21.0} để
  theo dõi phiên bản image của các container được chỉ định.
  \item \textbf{restart:} Thiết lập \texttt{unless-stopped} để dịch vụ tự khởi
  động lại khi Docker daemon hoặc máy chủ khởi động lại, trừ khi dịch vụ được
  dừng chủ động.
  \item \textbf{\texttt{/var/run/docker.sock:/var/run/docker.sock}:} Gắn Docker
  socket của máy Host để Watchtower có thể kiểm tra image, dừng và tạo lại
  container khi có bản cập nhật.
  \item \textbf{\texttt{/root/.docker/config.json:/config.json:ro}:} Gắn thông
  tin xác thực Docker ở chế độ chỉ đọc, cho phép Watchtower truy cập image
  private từ GitHub Container Registry.
  \item \textbf{\texttt{-{}-interval 60}:} Thực hiện kiểm tra registry theo chu
  kỳ 60 giây để phát hiện image tag \texttt{latest} mới.
  \item \textbf{\texttt{-{}-cleanup}:} Tự động dọn dẹp các image cũ sau khi
  container được cập nhật, giúp tiết kiệm dung lượng lưu trữ.
  \item \textbf{\texttt{-{}-debug}:} Bật log chi tiết để hỗ trợ theo dõi tiến
  trình và gỡ lỗi quá trình cập nhật.
  \item \textbf{vproapp:} Giới hạn phạm vi theo dõi cho container vproapp, tránh
  cập nhật ngoài ý muốn đối với các dịch vụ hạ tầng khác.
\end{itemize}

\ctufigure{image17.png}{0.94}{Dịch vụ Watchtower}{fig:watchtower-service}

\subsubsection{Cấu hình dịch vụ Ứng dụng Backend}
\textbf{Dịch vụ Ứng dụng Backend được cấu hình như sau:}
\begin{itemize}
  \item \textbf{image:} Kéo image ứng dụng trực tiếp từ GitHub Container
  Registry (GHCR) (\texttt{ghcr.io/duytechie/actionprofile:latest}).
  \item \textbf{labels:} Đánh dấu nhãn \texttt{enable: "true"} để dịch vụ
  Watchtower nhận diện và theo dõi việc cập nhật image mới.
  \item \textbf{depends\_on:} Chỉ định thứ tự phụ thuộc. vproapp sẽ chỉ được
  khởi động sau khi 3 dịch vụ nền tảng (vprodb, vprocache01, vpromq01) đã được
  khởi chạy.
\end{itemize}

\ctufigure{image18.png}{0.80}{Dịch vụ Ứng dụng Backend}{fig:vproapp}

\subsubsection{Cấu hình dịch vụ Web Proxy vproweb}
\textbf{Dịch vụ Web Proxy vproweb được cấu hình như sau:}
\begin{itemize}
  \item \textbf{build.context:} Xây dựng image từ thư mục
  \nolinkurl{./Docker-files/web}. Dockerfile sử dụng image Nginx làm nền tảng, xoá
  cấu hình mặc định và sao chép file \texttt{nginvproapp.conf} vào thư mục cấu
  hình của Nginx.
  \item \textbf{image:} Đặt tên image sau khi build là
  \nolinkurl{vprocontainers/vprofileweb} để thuận tiện cho việc quản lý và triển
  khai lại dịch vụ.
  \item \textbf{ports (80:80):} Công khai cổng HTTP 80 của container ra máy Host,
  tạo điểm truy cập web cho người dùng.
  \item \textbf{nginvproapp.conf:} Khai báo upstream vproapp trỏ đến
  \texttt{vproapp:8080} và sử dụng \texttt{proxy\_pass} để chuyển tiếp mọi yêu
  cầu từ Nginx đến ứng dụng Backend thông qua mạng nội bộ Docker.
  \item \textbf{depends\_on:} Chỉ định vproweb khởi động sau vproapp, bảo đảm
  lớp proxy được triển khai sau dịch vụ ứng dụng Backend.
\end{itemize}

\ctufigure{image19.png}{0.66}{Dịch vụ Web Proxy vproweb}{fig:vproweb}

\subsubsection{Theo dõi quá trình triển khai tự động}
Sau khi image ứng dụng được phát hành, quá trình triển khai có thể được theo
dõi qua kết quả workflow và log của Watchtower. Log cần thể hiện các bước phát
hiện image mới, dừng container vproapp và khởi tạo lại container với cấu hình
đã cập nhật.

\ctufigure{image20.png}{1.0}{Kết quả thực thi workflow triển khai ứng dụng}{fig:workflow-result}

\ctufigure{image21.png}{0.62}{Kết quả Workflow trong VSCode}{fig:workflow-vscode}

Log của Watchtower được sử dụng để xác nhận hệ thống đã kéo image mới và cập
nhật container vproapp thành công.

\ctufigure{image22.png}{1.0}{Log của Watchtower khi cập nhật container vproapp}{fig:watchtower-log-deploy}

% -----------------------------------------------------------------------------
% CHƯƠNG IV
% -----------------------------------------------------------------------------
\chapter{ĐÁNH GIÁ KIỂM THỬ}
\label{chap:evaluation}
\section{Mục tiêu kiểm thử}
Kiểm thử quy trình CI/CD đã xây dựng cho ứng dụng VProfile để xác nhận một thay
đổi mã nguồn có được tích hợp, kiểm tra, đóng gói, phát hành và triển khai đúng
theo luồng tự động hay không.
\begin{itemize}
  \item \textbf{Phát hiện lỗi cấu hình:} Kiểm tra các lỗi có thể làm workflow
  không được kích hoạt, Maven không tạo được artifact, image không được phát
  hành hoặc container không được cập nhật.
  \item \textbf{Đối chiếu yêu cầu:} Đánh giá các kết quả thực tế với yêu cầu đặt
  ra cho pipeline: build một lần, quét bảo mật, build đa kiến trúc, công bố
  image lên GHCR và tự động cập nhật dịch vụ vproapp.
  \item \textbf{Quan sát log file:} Ghi nhận log Maven, kết quả GitHub Actions,
  artifact, manifest image và log Watchtower để phục vụ việc đánh giá, bảo trì
  và khắc phục sự cố về sau.
\end{itemize}

\section{Phạm vi kiểm thử}
Việc kiểm thử được thực hiện theo các nhóm sau:
\begin{itemize}
  \item \textbf{Kiểm thử cấu hình và kích hoạt Workflow:} Xác nhận workflow được
  kích hoạt khi có push hoặc pull request vào nhánh main, đồng thời phân quyền
  đọc mã nguồn và ghi package được khai báo đúng.
  \item \textbf{Kiểm thử tích hợp liên tục:} Xác nhận Maven có thể biên dịch,
  chạy vòng đời test, đóng gói WAR và tải artifact vprofile-app lên GitHub
  Actions.
  \item \textbf{Kiểm thử chất lượng và bảo mật:} Xác nhận báo cáo Checkstyle và
  file \texttt{trivy-results.json} được sinh ra, có thể tải lại để đánh giá mã
  nguồn và thư viện phụ thuộc.
  \item \textbf{Kiểm thử phát hành và triển khai:} Xác nhận image AMD64 và ARM64
  được gộp thành manifest trên GHCR; Watchtower phát hiện image latest mới và
  cập nhật đúng container vproapp.
\end{itemize}

Ngoài phạm vi kiểm thử:
\begin{itemize}
  \item Kiểm thử khả năng chịu tải, kiểm thử hiệu năng chuyên sâu, kiểm thử xâm
  nhập và kiểm thử đầy đủ mọi chức năng giao diện người dùng chưa thuộc phạm vi
  của niên luận này.
  \item Workflow này chỉ thực hiện các bước test cơ bản. Chưa thực hiện các quy
  trình test và security test đầy đủ và chuyên nghiệp.
\end{itemize}

\section{Kế hoạch kiểm thử}
Các nội dung được kiểm thử:
\begin{itemize}
  \item \textbf{Kích hoạt pipeline:} Sự kiện push và pull request vào nhánh
  main.
  \item \textbf{Build và artifact:} Lệnh \texttt{mvn install}, file WAR tại
  \texttt{target} và artifact vprofile-app.
  \item \textbf{Quét và báo cáo:} Job Security Scan, file
  \texttt{trivy-results.json} và báo cáo Checkstyle.
  \item \textbf{Image đa kiến trúc:} Job BUILD\_AMD64, BUILD\_ARM64 và
  PUBLISH\_TO\_GHCR cùng các tag commit, latest trên GHCR.
  \item \textbf{Triển khai tự động:} Chuỗi log Watchtower từ lúc phát hiện image
  mới đến khi khởi tạo lại vproapp.
\end{itemize}

\section{Các bước kiểm thử và kết quả ghi nhận}
Lúc này server Ubuntu để dựng trang web Java đã được khởi động sẵn.

Mã nguồn được thay đổi và đẩy lên GitHub đồng thời kích hoạt GitHub Workflow.

\ctufigure{image23.jpeg}{1.0}{Git commit}{fig:git-commit}

\subsection{Kiểm thử workflow và artifact}
GitHub Workflow được kích hoạt. Kết quả mong đợi là các job hiển thị Success,
\textbf{artifact vprofile-app} và \textbf{trivy-scan-results} có thể tải về, còn
manifest đa kiến trúc được tạo sau khi hai image theo digest hoàn tất.

\ctufigure{image24.png}{1.0}{Summary của GitHub Actions và danh sách Artifacts}{fig:gha-summary}

\subsection{Kiểm thử build và test Maven}
Kịch bản chạy lệnh \texttt{mvn test} nhằm xác nhận dự án có thể biên dịch và
các test case được Maven nhận diện. Kết quả ghi nhận BUILD SUCCESS nhưng
Tests run: 0, Failures: 0, Errors: 0.

\ctufigure{image25.png}{1.0}{Kết quả mvn test}{fig:mvn-test}

\subsection{Kiểm thử phát hành image đa kiến trúc}
Kịch bản kiểm tra package trên GHCR xác nhận cùng một tag \texttt{latest} có
thể tham chiếu đến image \texttt{linux/amd64} và \texttt{linux/arm64}.

\ctufigure{image26.png}{1.0}{Package trên GHCR}{fig:ghcr-package}

\subsection{Kiểm thử cập nhật tự động bằng Watchtower}
Sau đó nếu workflow chạy thành công thì image latest mới đã xuất hiện trên
GHCR. Kết quả mong đợi là Watchtower lần lượt kéo image, phát hiện image mới,
dừng vproapp cũ và khởi động lại vproapp với cấu hình được giữ nguyên. Các dịch
vụ hạ tầng khác không nằm trong phạm vi cập nhật của Watchtower.

Thực hiện kết nối SSH vào server Ubuntu.

\ctufigure{image27.png}{0.93}{vagrant ssh}{fig:vagrant-ssh}

Kiểm tra log của watchtower bằng lệnh sau:

\begin{lstlisting}[language=bash]
docker compose logs -f watchtower | rg --line-buffered 'Image pull completed|Found new image|Stopping container|Starting container with updated configuration|Update session completed'
\end{lstlisting}

\ctufigure{image22.png}{1.0}{Log Watchtower}{fig:watchtower-log-test}

Sau đó người dùng sẽ nhận được bản cập nhật mới một cách tự động sau lần
refresh lại trang web.

\ctufigure{image28.png}{0.56}{Phiên bản web mới được cập nhật}{fig:web-updated}

\subsection{Đánh giá kiểm thử}
Quy trình đã đáp ứng được mục tiêu tự động hóa việc build, test, security test,
phát hành image đa kiến trúc và triển khai lại dịch vụ của ứng dụng web.

\clearpage
\phantomsection
\addcontentsline{toc}{chapter}{TÀI LIỆU THAM KHẢO}
\begin{thebibliography}{9}
  \bibitem{vprofile}
  vprofile-project, \url{https://github.com/hkhcoder/vprofile-project}.

  \bibitem{200lab-cicd}
  200Lab, ``CI/CD là gì?'', \url{https://200lab.io/blog/ci-cd-la-gi}.

  \bibitem{push-ghcr}
  GitHub Marketplace, ``Push to GHCR'',
  \url{https://github.com/marketplace/actions/push-to-ghcr}.

  \bibitem{login-action}
  Docker, ``Login Action'', \url{https://github.com/docker/login-action}.

  \bibitem{blackduck-cicd}
  Black Duck, ``Continuous Integration and Continuous Delivery/Deployment'',
  \url{https://www.blackduck.com/glossary/what-is-cicd.html}.
\end{thebibliography}

\end{document}
